\section{Descripcion de la solucion}

\subsection{Arquitectura}

La solucion esta dividida en modulos. \texttt{common/} contiene el protocolo,
wrappers de sockets y hashing. \texttt{server/} implementa el servidor central
y su registro de pares. \texttt{client/} contiene el arranque, escaneo de
carpetas, API hacia el servidor y REPL. \texttt{transfer/} implementa el
listener del puerto de datos, el envio de rangos y el receptor que ensambla
segmentos. \texttt{search/} contiene vecinos, flooding y agregacion de
respuestas.

\subsection{Servidor central}

El servidor escucha conexiones TCP y crea un hilo por cliente aceptado. Cada
mensaje debe tener version valida, opcode conocido y longitud de payload
correcta. El opcode \texttt{REGISTER} registra la IP observada en la conexion,
el puerto de datos anunciado y la lista de archivos del cliente. La respuesta
incluye hasta \texttt{P2P\_MAX\_NEIGHBORS} pares recientes, excluyendo al par
que acaba de registrarse. El opcode \texttt{FIND} busca coincidencias parciales
por nombre y retorna metadatos \texttt{S}, \texttt{H}, IP, puerto y nombre.
Los errores de version, formato u opcode se responden con \texttt{ERROR}.

\subsection{Cliente y comandos}

Al iniciar, el cliente recibe IP y puerto del servidor, puerto de datos y
carpeta compartida. Escanea recursivamente la carpeta, calcula hash y tamano,
registra los metadatos en el servidor, inicializa la busqueda distribuida y
arranca un listener unico para mensajes entre pares.

El REPL soporta:
\begin{itemize}
  \item \texttt{find -s nombre}: busqueda centralizada.
  \item \texttt{find -d nombre}: busqueda distribuida.
  \item \texttt{find nombre}: servidor primero; si falla o retorna cero
  resultados, usa busqueda distribuida.
  \item \texttt{request S H}: descarga usando los pares de la ultima busqueda
  mostrada.
\end{itemize}

\subsection{Transferencia segmentada}

El receptor filtra los resultados cacheados por \texttt{S} y \texttt{H},
elimina pares duplicados y divide el archivo entre los pares disponibles. Cada
hilo de segmento abre una conexion TCP al puerto de datos del par, envia
\texttt{TRANSFER\_REQ} con \texttt{hash}, \texttt{size}, \texttt{byte\_start}
y \texttt{byte\_end}, y recibe frames \texttt{TRANSFER\_DATA}. Los bytes se
escriben con \texttt{pwrite} en un archivo parcial. Si todos los segmentos
terminan bien, el archivo parcial se renombra a
\texttt{download\_S\_H.bin}.

\subsection{Listener unico del puerto de datos}

Para evitar que transferencia y busqueda distribuida compitan por el mismo
puerto, el cliente usa un solo listener. El dispatcher lee el header del frame
y enruta por opcode: \texttt{TRANSFER\_REQ} hacia el sender,
\texttt{QUERY\_FLOOD} hacia la logica de flooding,
\texttt{QUERY\_RESULT} hacia el agregador, y opcodes invalidos hacia
\texttt{ERROR}.

\subsection{Busqueda distribuida}

La lista de vecinos se siembra con la respuesta de registro del servidor y se
actualiza con resultados observados. \texttt{find -d} genera un ID de consulta
de 64 bits en hexadecimal, marca el ID como visto localmente, envia
\texttt{QUERY\_FLOOD} a los vecinos y abre una ventana de agregacion. Cada par
que recibe una consulta nueva escanea su carpeta compartida en ese momento; por
eso tambien puede encontrar archivos agregados despues del registro central.
Las respuestas \texttt{QUERY\_RESULT} se envian directamente al originador.

Se mantuvo congelado \texttt{common/protocol.h}. Como el mensaje congelado
incluye el originador pero no el puerto de datos del remitente inmediato, el
control de ciclos depende de la cache de IDs vistos. En las pruebas locales,
esto permite trabajar con varios clientes en la misma maquina sin confundir
pares distintos que comparten IP.

\subsection{Tolerancia a hot-unplug}

Antes de escanear o abrir archivos se revisa acceso a la carpeta. Si la carpeta
compartida desaparece, el cliente registra una advertencia y continua vivo. Las
fallas de lectura se propagan como errores de transferencia o como busquedas
locales sin resultados, pero no terminan el servidor ni el proceso cliente.
